\section{Related work}
\label{sec:related}

\paragraph{Grounding language models with structure.} Retrieval-augmented
generation established the general pattern of conditioning generation on
retrieved evidence \cite{lewis2020rag}, and a substantial line of work has since
argued that \emph{structured} evidence outperforms passage retrieval when
questions require aggregation or multi-hop traversal, most visibly the
community around graph-structured retrieval \cite{edge2024graphrag} and hybrid
graph/vector systems \cite{sarmah2024hybridrag}. Surveys of knowledge-graph
augmentation of language models \cite{pan2024unifying} map the design space.
Our results are consistent with that finding and refine it. In an industrial
setting the decisive structure is not the entity graph that general-domain
GraphRAG systems extract from text; it is the \emph{engineering} structure that
already exists in standards: units, ranges, access rights, topology, state
machines and interlocks. Text-derived graphs cannot recover this material
because it was never written in text. Two consequences follow: industrial knowledge-graph
construction should start from the control and engineering systems rather than
from document extraction, and the corresponding benchmarks need to include
actions and constraints, not only questions.

\paragraph{Measuring truthfulness and faithfulness.} Benchmarks such as
TruthfulQA \cite{lin2022truthfulqa} and fine-grained factuality measures such as
FActScore \cite{min2023factscore} evaluate free-text claims against reference
knowledge. Industrial answers differ in kind: they are numeric, unit-bearing and
tied to a specific asset at a specific time, which makes faithfulness
\emph{decidable} rather than estimable. Our \textsc{f-value} check
(\S\ref{sec:layer3}) is one instance. An answer modelled as an observation can
be compared exactly, by query, with the observation it claims to report.

\paragraph{Tool use and function calling.} The agentic turn rests on models
emitting structured calls against declared tool schemas
\cite{schick2023toolformer,yao2023react}, and the evaluation of that capability
has become a benchmark discipline in its own right \cite{patil2024gorilla}.
Those evaluations hold the schema fixed and vary the model. We do the reverse:
Experiment~E5 holds the model out entirely and varies the \emph{schema},
measuring how much of the error surface a semantically derived schema removes
before the model is even invoked.

\paragraph{Constraints and neuro-symbolic verification.} The idea of checking
generated output against a formal specification recurs across program synthesis,
constrained decoding and schema-guided generation. SHACL
\cite{knublauch2017shacl} provides the industrial-strength instance of this for
graph data, and the shapes literature has explored validation and repair. What
is new here is the application to \emph{agent actions in a physical plant}, where
the constraint set is not a data-quality convention but the plant's operating
envelope, and where we can therefore quantify detection and false alarms per
fault family.

\paragraph{Industrial semantic interoperability.} There is a long tradition of
work on \opcua{} information modelling, on the Asset Administration Shell as an
Industrie~4.0 digital twin \cite{idta2023aas}, on AutomationML for engineering
exchange, and on ontology-based integration in manufacturing. Reviews in this
tradition evaluate against interoperability criteria such as coverage, mapping
effort and conformance. That is the right frame for machine-to-machine
integration and the wrong one for agents. The agent-usability framework of
\S\ref{sec:framework} is our attempt to supply the missing frame.

\paragraph{Domain ontologies and units.} SOSA/SSN \cite{haller2019ssn} for
observations, Brick \cite{balaji2016brick} for building equipment, SAREF for
smart devices, and QUDT \cite{qudt} for quantities and units are the vocabularies
this paper leans on most. Their design has been evaluated for modelling adequacy;
their value \emph{as agent infrastructure}, and in QUDT's case its
disproportionate value, has not, to our knowledge, been quantified before.

\paragraph{LLMs in automation engineering.} A growing body of work applies
language models to control-code generation, engineering document processing and
operator assistance, generally reporting that domain context materially improves
results and that verification remains the open problem. Our study is
complementary in method: rather than measuring one pipeline end-to-end, we
isolate the representational variable and bound what any pipeline built on it can
achieve.

\paragraph{Positioning.} To our knowledge this is the first work to (i) assess
the industrial semantic technology landscape explicitly against agent-usability
criteria, (ii) ablate the semantic stack in a controlled testbed with a
model-independent accuracy ceiling, and (iii) quantify per-fault-family guardrail
efficacy and action-space fidelity for industrial agent actions.
Table~\ref{tab:novelty} places the paper against the closest lines of work on the
dimensions that distinguish it.

\begin{table*}[t]
\centering
\caption{Positioning against the closest related work. \checkmark{} denotes a
central contribution, (\checkmark) a partial or incidental one, and a blank its
absence. The combination in the final row, not any single column, is what we
claim is new.}
\label{tab:novelty}
\footnotesize
\setlength{\tabcolsep}{4pt}
\begin{tabular}{@{}L{4.3cm}cccccc@{}}
\toprule
\textbf{Work / line} & \textbf{Action}      & \textbf{Validator}   & \textbf{Naming} & \textbf{Model-indep.} & \textbf{Reprod.} & \textbf{Engineering} \\
                     & \textbf{legality}    & \textbf{false alarms} & \textbf{robustness} & \textbf{ceiling}   & \textbf{artefact} & \textbf{KG source} \\
\midrule
GraphRAG \cite{edge2024graphrag}           &            &            &            &            & (\checkmark) &            \\
HybridRAG \cite{sarmah2024hybridrag}       &            &            &            &            & (\checkmark) &            \\
KG-for-LLM survey \cite{pan2024unifying}   &            &            &            &            &            & (\checkmark) \\
Tool/function benchmarks \cite{patil2024gorilla} & (\checkmark) &      &            &            & \checkmark &            \\
SHACL validation \cite{knublauch2017shacl} &            & (\checkmark) &          &            & \checkmark &            \\
Semantic interoperability reviews \cite{idta2023aas} & &          &            &            &            & (\checkmark) \\
\midrule
\textbf{This work}                         & \checkmark & \checkmark & \checkmark & \checkmark & \checkmark & \checkmark \\
\bottomrule
\end{tabular}
\end{table*}

\paragraph{What we do not claim.} To keep the contribution precise: we do
\emph{not} propose a new ontology or vocabulary, and we reuse QUDT, SOSA/SSN,
Brick and PROV-O as published; we do \emph{not} claim SHACL is the only, or the
optimal, constraint mechanism, and \S\ref{sec:results} reports a non-RDF
procedural validator as a baseline against exactly this objection; we do
\emph{not} claim the agent-usability scores of Table~\ref{tab:auf-matrix} are the
only defensible assessment, only that they are an explicit and reproducible one;
and we do \emph{not} claim a specific model's accuracy, which is why the central
measure is a representation ceiling rather than a leaderboard number.
